% ============================================================================
% D1 2026-09-29 (change 7): where the headline's P&L comes from.  A subsection that
% continues Appendix D (input from main.tex right after sections/appendix_running).
% Every number is a \cm... macro and the table is generated/table_close_pnl.tex, both
% written by writeup/make_table_close_main.py from
% results/close_pnl_decomp/research_scorer/*.csv (experiments/close_pnl_decomposition.py
% re-run under the 16:00-bar recalibration of Table 7).  The same decomposition under the
% session-bar recalibration is results/close_pnl_decomp/SUMMARY.md; it is not quoted here.
% ============================================================================
\ifdefined\cmNdays\else% AUTO-GENERATED by writeup/make_table_close_main.py -- do not edit.
% Numbers quoted in the prose of writeup/sections/results_close_option.tex; each is read
% from results/close_master_table/master_table.csv (16:00-bar recalibration).
\newcommand{\cmNdays}{866}
\newcommand{\cmNrows}{28}
\newcommand{\cmShortMean}{0.014}
\newcommand{\cmShortShMid}{0.20}
\newcommand{\cmShortShX}{\ensuremath{-}0.27}
\newcommand{\cmShortShCiMid}{[\ensuremath{-}0.63, 1.14]}
\newcommand{\cmShortShCiX}{[\ensuremath{-}1.10, 0.63]}
\newcommand{\cmPaperShMin}{1.00}
\newcommand{\cmPaperShMax}{1.54}
\newcommand{\cmPaperXMin}{0.53}
\newcommand{\cmPaperXMax}{1.07}
\newcommand{\cmPaperVsShortMin}{+0.80}
\newcommand{\cmPaperVsShortMax}{+1.33}
\newcommand{\cmPaperVsShortAboveN}{one}
\newcommand{\cmPaperVsShortAboveWho}{LightGBM}
\newcommand{\cmPaperVsShortAboveCi}{+1.27 [+0.05, +2.41]}
\newcommand{\cmPaperSameMin}{79}
\newcommand{\cmSellPaperMin}{62}
\newcommand{\cmSellPaperMax}{73}
\newcommand{\cmSellHead}{60}
\newcommand{\cmPaperSameMax}{96}
\newcommand{\cmRefQ}{0.1058}
\newcommand{\cmRefShMid}{1.00}
\newcommand{\cmRefShX}{0.53}
\newcommand{\cmRefBuy}{36}
\newcommand{\cmRefVsShort}{+0.80 [\ensuremath{-}0.59, +2.01]}
\newcommand{\cmFomcSh}{1.12}
\newcommand{\cmFomcD}{+0.12 [\ensuremath{-}0.35, +0.63]}
\newcommand{\cmLassoFShMid}{1.54}
\newcommand{\cmLassoFVsRef}{+0.54 [+0.07, +1.06]}
\newcommand{\cmLassoFSame}{96}
\newcommand{\cmNlFamN}{7}
\newcommand{\cmNlMasterN}{147}
\newcommand{\cmPoolShMin}{0.99}
\newcommand{\cmPoolShMax}{1.24}
\newcommand{\cmSubShMin}{1.22}
\newcommand{\cmSubShMax}{1.90}
\newcommand{\cmNlShMin}{0.87}
\newcommand{\cmNlShMax}{1.63}
\newcommand{\cmTabVsRefN}{27}
\newcommand{\cmTabVsRefAbove}{2}
\newcommand{\cmTabVsRefBelow}{none}
\newcommand{\cmMasterN}{220}
\newcommand{\cmMasterVsRefN}{219}
\newcommand{\cmMasterVsRefAbove}{10}
\newcommand{\cmHQ}{0.1005}
\newcommand{\cmHQPct}{\ensuremath{-}5.0}
\newcommand{\cmHQCi}{[\ensuremath{-}0.0113, +0.0005]}
\newcommand{\cmHDM}{\ensuremath{-}1.42}
\newcommand{\cmHDMp}{0.156}
\newcommand{\cmHShMid}{1.90}
\newcommand{\cmHShX}{1.44}
\newcommand{\cmHMean}{0.134}
\newcommand{\cmHHit}{54.6}
\newcommand{\cmHBuy}{40.2}
\newcommand{\cmHVsRefMid}{+0.90 [+0.19, +1.66]}
\newcommand{\cmHVsShortMid}{+1.70 [+0.09, +3.11]}
\newcommand{\cmHVsRefX}{+0.91 [+0.19, +1.67]}
\newcommand{\cmHVsShortX}{+1.71 [+0.12, +3.12]}
\newcommand{\cmHRankSh}{2}
\newcommand{\cmHRankQ}{63}
\newcommand{\cmTopShLabel}{per-bar lasso}
\newcommand{\cmTopShInputs}{\texttt{live\_vix\_only}}
\newcommand{\cmTopShMid}{1.91}
\newcommand{\cmVsHLasso}{\ensuremath{-}0.08 [\ensuremath{-}0.89, +0.79]}
\newcommand{\cmVsHLassoSame}{89}
\newcommand{\cmVsHLassoQPct}{\ensuremath{-}4.1}
\newcommand{\cmVsHLassoDM}{\ensuremath{-}1.20}
\newcommand{\cmVsHEnet}{\ensuremath{-}0.49 [\ensuremath{-}1.08, +0.14]}
\newcommand{\cmVsHEnetSame}{91}
\newcommand{\cmVsHEnetQPct}{\ensuremath{-}3.7}
\newcommand{\cmVsHEnetDM}{\ensuremath{-}1.38}
\newcommand{\cmVsHRidgeAll}{\ensuremath{-}0.24 [\ensuremath{-}1.07, +0.54]}
\newcommand{\cmVsHRidgeAllSame}{86}
\newcommand{\cmVsHRidgeAllQPct}{\ensuremath{-}0.2}
\newcommand{\cmVsHRidgeAllDM}{\ensuremath{-}0.06}
\newcommand{\cmVsHTop}{+0.01 [\ensuremath{-}0.76, +0.81]}
\newcommand{\cmVsHTopSame}{88}
\newcommand{\cmVsHTopQPct}{\ensuremath{-}3.6}
\newcommand{\cmVsHTopDM}{\ensuremath{-}1.04}
\newcommand{\cmVsHLassoSh}{1.83}
\newcommand{\cmVsHN}{219}
\newcommand{\cmVsHAbove}{none}
\newcommand{\cmVsHBelow}{43}
\newcommand{\cmXLossMin}{0.46}
\newcommand{\cmXLossMax}{0.48}
\newcommand{\cmLdSpec}{\ensuremath{-}0.01 [\ensuremath{-}0.73, +0.76]}
\newcommand{\cmLdSpecQ}{\ensuremath{-}1.5 [\ensuremath{-}5.8, +2.3]}
\newcommand{\cmLdSpecDM}{\ensuremath{-}0.75}
\newcommand{\cmLdPerBarBase}{\ensuremath{-}0.03 [\ensuremath{-}0.66, +0.64]}
\newcommand{\cmLdPerBarBaseQ}{\ensuremath{-}11.2 [\ensuremath{-}15.7, \ensuremath{-}6.9]}
\newcommand{\cmLdPerBarBaseDM}{\ensuremath{-}4.59}
\newcommand{\cmLdPerBarAll}{+0.68 [\ensuremath{-}0.12, +1.45]}
\newcommand{\cmLdPerBarAllQ}{\ensuremath{-}3.6 [\ensuremath{-}10.7, +3.1]}
\newcommand{\cmLdPerBarAllDM}{\ensuremath{-}0.92}
\newcommand{\cmLdPerBarLive}{+0.78 [\ensuremath{-}0.08, +1.52]}
\newcommand{\cmLdPerBarLiveQ}{\ensuremath{-}2.8 [\ensuremath{-}7.9, +2.6]}
\newcommand{\cmLdPerBarLiveDM}{\ensuremath{-}0.87}
\newcommand{\cmLdInAll}{+0.24 [\ensuremath{-}0.54, +1.07]}
\newcommand{\cmLdInAllQ}{+0.2 [\ensuremath{-}4.9, +5.5]}
\newcommand{\cmLdInAllDM}{+0.06}
\newcommand{\cmLdInBase}{+0.69 [\ensuremath{-}0.23, +1.70]}
\newcommand{\cmLdInBaseQ}{+2.4 [\ensuremath{-}5.8, +12.3]}
\newcommand{\cmLdInBaseDM}{+0.53}
\newcommand{\cmLdEstMin}{\ensuremath{-}0.63}
\newcommand{\cmLdEstMax}{\ensuremath{-}0.08}
\newcommand{\cmLdN}{26}
\newcommand{\cmLdNPath}{14}
\newcommand{\cmLdNCampaign}{12}
\newcommand{\cmLdResolved}{none}
\newcommand{\cmLdCDedup}{+0.00 [+0.00, +0.00]}
\newcommand{\cmLdCDedupQ}{+0.0 [+0.0, +0.0]}
\newcommand{\cmLdCMaskTten}{\ensuremath{-}0.17 [\ensuremath{-}0.80, +0.39]}
\newcommand{\cmLdCMaskTtenQ}{+1.0 [\ensuremath{-}0.4, +2.3]}
\newcommand{\cmLdCMaskTtenDM}{+1.31}
\newcommand{\cmLdCMaskTone}{\ensuremath{-}0.06 [\ensuremath{-}0.49, +0.37]}
\newcommand{\cmLdCMaskToneQ}{\ensuremath{-}1.0 [\ensuremath{-}2.4, +0.3]}
\newcommand{\cmLdCMaskToneDM}{\ensuremath{-}1.47}
\newcommand{\cmLdCMaskRsten}{\ensuremath{-}0.12 [\ensuremath{-}0.67, +0.43]}
\newcommand{\cmLdCMaskRstenQ}{\ensuremath{-}0.9 [\ensuremath{-}2.7, +0.8]}
\newcommand{\cmLdCMaskRstenDM}{\ensuremath{-}0.93}
\newcommand{\cmLdCMaskRsone}{+0.01 [\ensuremath{-}0.55, +0.58]}
\newcommand{\cmLdCMaskRsoneQ}{\ensuremath{-}0.1 [\ensuremath{-}2.1, +1.8]}
\newcommand{\cmLdCMaskRsoneDM}{\ensuremath{-}0.10}
\newcommand{\cmLdCMaskLstm}{\ensuremath{-}0.05 [\ensuremath{-}0.69, +0.57]}
\newcommand{\cmLdCMaskLstmQ}{+0.3 [\ensuremath{-}3.0, +3.4]}
\newcommand{\cmLdCMaskLstmDM}{+0.15}
\newcommand{\cmLdCCadLgbm}{\ensuremath{-}0.17 [\ensuremath{-}1.11, +0.59]}
\newcommand{\cmLdCCadLgbmQ}{\ensuremath{-}3.6 [\ensuremath{-}7.4, \ensuremath{-}0.8]}
\newcommand{\cmLdCCadLgbmDM}{\ensuremath{-}2.73}
\newcommand{\cmLdCCadLstm}{\ensuremath{-}0.17 [\ensuremath{-}0.80, +0.48]}
\newcommand{\cmLdCCadLstmQ}{\ensuremath{-}2.5 [\ensuremath{-}7.8, +3.2]}
\newcommand{\cmLdCCadLstmDM}{\ensuremath{-}0.77}
\newcommand{\cmLdCTuneRs}{+0.21 [\ensuremath{-}0.61, +1.07]}
\newcommand{\cmLdCTuneRsQ}{+4.1 [\ensuremath{-}0.2, +8.9]}
\newcommand{\cmLdCTuneRsDM}{+1.70}
\newcommand{\cmLdCTuneOp}{\ensuremath{-}0.55 [\ensuremath{-}1.37, +0.24]}
\newcommand{\cmLdCTuneOpQ}{+7.4 [+1.4, +16.3]}
\newcommand{\cmLdCTuneOpDM}{+2.32}
\newcommand{\cmLdCTuneOpPer}{\ensuremath{-}0.55 [\ensuremath{-}1.42, +0.30]}
\newcommand{\cmLdCTuneOpPerQ}{\ensuremath{-}2.8 [\ensuremath{-}9.2, +2.9]}
\newcommand{\cmLdCTuneOpPerDM}{\ensuremath{-}0.98}
\newcommand{\cmLdCSeq}{+0.62 [\ensuremath{-}0.26, +1.54]}
\newcommand{\cmLdCSeqQ}{\ensuremath{-}7.5 [\ensuremath{-}20.9, +4.7]}
\newcommand{\cmLdCSeqDM}{\ensuremath{-}1.25}
\newcommand{\cmPNewBase}{22}
\newcommand{\cmPOldBase}{34}
\newcommand{\cmPNewLive}{232}
\newcommand{\cmPOldLive}{244}
\newcommand{\cmPNewAll}{628}
\newcommand{\cmPOldAll}{640}
\newcommand{\cmNDup}{12}
\newcommand{\cmDedupN}{55}
\newcommand{\cmDedupQMax}{0.09}
\newcommand{\cmDedupPosMax}{two}
\newcommand{\cmKeptBase}{17}
\newcommand{\cmKeptBaseMin}{17}
\newcommand{\cmKeptBaseMax}{17}
\newcommand{\cmKeptLive}{138}
\newcommand{\cmKeptLiveMin}{136}
\newcommand{\cmKeptLiveMax}{145}
\newcommand{\cmKeptAll}{369}
\newcommand{\cmKeptAllMin}{360}
\newcommand{\cmKeptAllMax}{378}
\newcommand{\cmMaskN}{36}
\newcommand{\cmMaskQMed}{0.7}
\newcommand{\cmMaskQAbove}{five}
\newcommand{\cmMaskShBelow}{one}
\newcommand{\cmTrVsRidgeN}{72}
\newcommand{\cmTrVsRidgeQAbove}{32}
\newcommand{\cmTrVsRidgeShBelow}{two}
\newcommand{\cmCadN}{nine}
\newcommand{\cmCadQBelow}{three}
\newcommand{\cmCadShUp}{+0.52 [+0.10, +1.02]}
\newcommand{\cmRsQAbove}{6}
\newcommand{\cmRsN}{9}
\newcommand{\cmOpN}{36}
\newcommand{\cmOpQAbove}{19}
\newcommand{\cmTpN}{27}
\newcommand{\cmTpQBelow}{8}
\newcommand{\cmBudN}{36}
\newcommand{\cmBudQAbove}{one}
\newcommand{\cmBudLastMin}{27}
\newcommand{\cmBudLastMax}{40}
\newcommand{\cmBudLastExch}{20}
\newcommand{\cmLstmCadBase}{\ensuremath{-}1.9}
\newcommand{\cmLstmCadLive}{\ensuremath{-}2.5}
\newcommand{\cmLstmCadAll}{\ensuremath{-}10.1}
\newcommand{\cmLstmCadBelow}{two}
\newcommand{\cmLstmVsRidgeMin}{+6.8}
\newcommand{\cmLstmVsRidgeMax}{+30.0}
\newcommand{\cmIntraVsRidgeMin}{+9.2}
\newcommand{\cmIntraVsRidgeMax}{+11.5}
\newcommand{\cmIntraAllQ}{\ensuremath{-}14.3 [\ensuremath{-}26.9, \ensuremath{-}2.0]}
\newcommand{\cmNlBestSh}{1.90}
\newcommand{\cmNlBestVsH}{\ensuremath{-}0.00 [\ensuremath{-}0.56, +0.59]}
\newcommand{\cmNlVsHAbove}{none}
\newcommand{\cmClassChunks}{541}
\newcommand{\cmClassPts}{20}
\newcommand{\cmClassLgbmPick}{18}
\newcommand{\cmClassMaxRelMin}{4.5}
\newcommand{\cmClassMaxRelMax}{5.5}
\newcommand{\cmPnlTotal}{116.2}
\newcommand{\cmPnlTotalX}{87.4}
\newcommand{\cmPnlTopTenShare}{42}
\newcommand{\cmPnlTopTwentyShare}{68}
\newcommand{\cmPnlTopTwentyShareX}{86}
\newcommand{\cmPnlExTwenty}{+37.6}
\newcommand{\cmPnlExTwentySh}{0.76}
\newcommand{\cmPnlMktTwenty}{14}
\newcommand{\cmPnlMktTwentyExp}{8.0}
\newcommand{\cmPnlMktTwentyP}{0.006}
\newcommand{\cmPnlDaysAll}{40}
\newcommand{\cmPnlDaysAllPct}{4.6}
\newcommand{\cmPnlMEn}{52}
\newcommand{\cmPnlMEbuys}{12}
\newcommand{\cmPnlME}{\ensuremath{-}6.0 [\ensuremath{-}26.1, +12.7]}
\newcommand{\cmPnlMEShort}{\ensuremath{-}24.5 [\ensuremath{-}43.5, \ensuremath{-}7.3]}
\newcommand{\cmPnlNotMEn}{735}
\newcommand{\cmPnlNotME}{+116.9 [+59.8, +174.1]}
\newcommand{\cmPnlDiffBuyShare}{100}
\newcommand{\cmPnlDiff}{+103.7 [+10.7, +200.1]}
\newcommand{\cmPnlDiffX}{+104.9 [+12.1, +201.8]}
\newcommand{\cmPnlDiffExTen}{+6.0 [\ensuremath{-}59.8, +72.7]}
\newcommand{\cmPnlDiffHighVix}{+2.3 [\ensuremath{-}39.0, +48.3]}
\newcommand{\cmPnlHighVixN}{395}
\newcommand{\cmSvN}{31}
\newcommand{\cmSvAbove}{none}
\newcommand{\cmSvBestMidSh}{1.94}
\newcommand{\cmSvBestMid}{+0.04 [\ensuremath{-}0.04, +0.14]}
\newcommand{\cmSvWingN}{6}
\newcommand{\cmSvWingMid}{2}
\newcommand{\cmSvWingX}{6}
\newcommand{\cmSvCostStraddle}{2.9}
\newcommand{\cmSvCostIbfMin}{4.0}
\newcommand{\cmSvCostIbfMax}{6.1}
\fi
\subsection{Where the Headline's P\&L Comes From}
\label{sec:app_close_pnl}

\paragraph{Question.} Which days carry the P\&L of the last-30-min trade under
the headline forecast of Section~\ref{sec:close_option_results}, and does the
rule's advantage over always short come from a calendar effect? The series are
the headline's $\mathrm{sign}(s)$ portfolio and always short, under the 16:00-bar
recalibration of Table~\ref{tab:close_option_rules}, on its \cmNdays{} days at
the midpoint. P\&L is in premium units: one unit of premium staked every day,
daily returns summed and never compounded. Intervals are 95\% circular block
bootstrap intervals (blocks of 21 days, $2{,}000$ draws) on the same resampled
days for every series, a cell's sum recomputed on the resampled days that fall
in it. The VIX terciles place the previous session's VIX close against the
terciles of every close since 1990 up to that session.

\begin{table}[H]
\centering
\caption{The headline's $\mathrm{sign}(s)$ P\&L by cell, in premium units at
the midpoint, against always short on the same days; the difference is the
headline minus always short, zero on every day the rule sells. Buys: the days
the rule is long.}
\label{tab:app_close_pnl}
\resizebox{\textwidth}{!}{% AUTO-GENERATED by writeup/make_table_close_main.py -- do not edit.
% Source: results/close_pnl_decomp/research_scorer/{calendar_cells,regime_cells,
% diff_attribution,tail_concentration}.csv (experiments/close_pnl_decomposition.py, the
% 16:00-bar recalibration; midpoint fills).
\begingroup\small\setlength{\tabcolsep}{4pt}
\begin{tabular}{lrrccc}
\toprule
cell & days & buys & $\mathrm{sign}(s)$, headline & always short & difference \\
\midrule
all days & 866 & 348 & +116.2 [+50.2, +183.0] & +12.5 [\ensuremath{-}41.5, +65.8] & +103.7 [+10.7, +200.1] \\
\addlinespace
month-end, last session & 52 & 12 & \ensuremath{-}6.0 [\ensuremath{-}26.1, +12.7] & \ensuremath{-}24.5 [\ensuremath{-}43.5, \ensuremath{-}7.3] & +18.5 [\ensuremath{-}1.2, +41.3] \\
month-end, the session before & 35 & 13 & \ensuremath{-}2.6 [\ensuremath{-}13.8, +7.8] & +6.7 [\ensuremath{-}4.3, +16.5] & \ensuremath{-}9.3 [\ensuremath{-}19.7, +1.2] \\
month-end, the session after & 44 & 19 & +8.0 [\ensuremath{-}4.6, +21.1] & +10.6 [\ensuremath{-}2.3, +22.4] & \ensuremath{-}2.6 [\ensuremath{-}21.4, +19.0] \\
more than one session from a month-end & 735 & 304 & +116.9 [+59.8, +174.1] & +19.7 [\ensuremath{-}27.3, +64.6] & +97.2 [+12.4, +184.1] \\
FOMC statement day & 33 & 12 & +1.3 [\ensuremath{-}9.9, +11.1] & \ensuremath{-}5.0 [\ensuremath{-}16.2, +5.2] & +6.3 [\ensuremath{-}2.2, +15.3] \\
\addlinespace
previous VIX close: low tercile & 172 & 105 & +15.4 [\ensuremath{-}23.1, +56.9] & \ensuremath{-}14.0 [\ensuremath{-}43.5, +12.3] & +29.4 [\ensuremath{-}20.7, +89.4] \\
previous VIX close: middle tercile & 299 & 120 & +77.1 [+36.0, +122.7] & +5.1 [\ensuremath{-}32.3, +44.2] & +72.0 [+7.4, +139.5] \\
previous VIX close: high tercile & 395 & 123 & +23.7 [\ensuremath{-}8.2, +57.8] & +21.4 [\ensuremath{-}10.2, +53.0] & +2.3 [\ensuremath{-}39.0, +48.3] \\
\addlinespace
2020 & 158 & 69 & +21.0 [\ensuremath{-}4.2, +51.1] & \ensuremath{-}4.8 [\ensuremath{-}25.2, +14.6] & +25.8 [\ensuremath{-}8.6, +64.4] \\
2021 & 158 & 54 & +20.4 [\ensuremath{-}9.4, +56.3] & \ensuremath{-}6.4 [\ensuremath{-}32.7, +17.7] & +26.8 [\ensuremath{-}19.4, +82.9] \\
2022 & 219 & 49 & +2.0 [\ensuremath{-}18.3, +23.7] & +22.3 [+0.6, +47.7] & \ensuremath{-}20.2 [\ensuremath{-}44.8, +2.5] \\
2023 & 248 & 124 & +47.0 [+6.9, +96.8] & +15.8 [\ensuremath{-}9.8, +44.0] & +31.3 [\ensuremath{-}15.5, +88.9] \\
2024 (to 30 April) & 83 & 52 & +25.8 [+0.3, +57.1] & \ensuremath{-}14.3 [\ensuremath{-}31.7, \ensuremath{-}0.5] & +40.1 [+5.9, +80.9] \\
\bottomrule
\end{tabular}
\endgroup
}
\end{table}

\paragraph{Finding.} The P\&L sits in a few long days. The headline makes
\cmPnlTotal{} premium units at the midpoint (\cmPnlTotalX{} crossed); its 10 % research_scorer/tail_concentration.csv
best days carry \cmPnlTopTenShare\% of that and its 20 best
\cmPnlTopTwentyShare\% (\cmPnlTopTwentyShareX\% crossed). All 20 best days are
buys and all 20 worst are sells; without the 20 best days the other days make
\cmPnlExTwenty{} (Sharpe \cmPnlExTwentySh), and \cmPnlDaysAll{} days --- % hit_payoff.csv days_for_100pct_of_pnl
\cmPnlDaysAllPct\% of the sample --- add up to the whole total. The rule owns
more of the straddle's largest payoffs than its buy rate gives: of the 20 days
the straddle paid most it bought \cmPnlMktTwenty{}, against % tail_concentration.csv mkt_top_k_bought
\cmPnlMktTwentyExp{} expected (hypergeometric $p=\cmPnlMktTwentyP$). Month-ends
are not the source: the long straddle pays on them on average (always short
makes \cmPnlMEShort{} on the \cmPnlMEn{} last sessions of a month), the rule % calendar_cells.csv
buys \cmPnlMEbuys{} of them and makes \cmPnlME{} there, against \cmPnlNotME{}
on the \cmPnlNotMEn{} days more than one session away. The difference to always
short, \cmPnlDiff{}, falls to \cmPnlDiffExTen{} without its 10 best days, and % diff_attribution.csv
in the high-VIX tercile the two make the same (\cmPnlDiffHighVix{} over % VIX tercile: high
\cmPnlHighVixN{} days).

\paragraph{Why it is not in the main text.} The main text keeps the three facts
that change its reading --- a tail trade, month-ends not the source, the
difference to always short and its interval; the cells are here.
